\appendix
\section{Checklist de reproducibilidad (plantilla)}
\begin{itemize}[leftmargin=*]
  \item \textbf{Datos}: fuente, versión, licencias, anonimización.
  \item \textbf{Código}: repositorio, commit hash, instrucciones de ejecución.
  \item \textbf{Entorno}: SO, versión de compiladores, dependencias, semillas.
  \item \textbf{Procedimiento}: pasos exactos para replicar resultados.
  \item \textbf{Resultados}: tablas/figuras generadas automáticamente en \texttt{build/}.
\end{itemize}

\section{Referencias Adicionales}
A continuación se presenta una tabla con referencias adicionales utilizadas en la investigación:

\begin{table}[htbp]
\centering
\tiny
\resizebox{0.8\textwidth}{!}{
\begin{tabular}{@{}p{0.15\textwidth}p{0.25\textwidth}p{0.5\textwidth}@{}}
\toprule
Autor / Fuente / Año & Título / Tema & Observaciones / Aplicaciones \\
\midrule
Barrios-Ulloa et al. (2025) & Agriculture 5.0 in Colombia: Opportunities Through the Emerging 6G Network & Analiza cómo tecnologías (IoT, IA, 6G) pueden modernizar la agricultura en Colombia: ideal para artículos sobre "Agricultura 5.0", conectividad rural, innovación tecnológica. MDPI \\
\midrule
Chowdhury et al. (2023) & The Role of Digital Agriculture in Transforming Rural Areas into Smart Villages & Describe el concepto "smart villages" y cómo la agricultura digital puede impulsar desarrollo rural sostenible — útil para visión de futuro, transformación rural, conectividad. arXiv \\
\midrule
Fonseca et al. (2019) & Agro 4.0: A Green Information System for Sustainable Agroecosystem Management & Presenta un sistema de información verde para gestión sostenible de agroecosistemas, útil para artículos sobre sostenibilidad, agroecología, uso de datos. arXiv \\
\midrule
Degada et al. (2021) & Smart Village: An IoT Based Digital Transformation & Examina la transformación digital de zonas rurales con IoT: buena base para abordar el uso de tecnología en zonas rurales, sensores, automatización. arXiv \\
\midrule
CEPAL (2020) & Digitalización y cambio: tecnologías digitales, agricultura y desarrollo rural & Análisis general del rol de TIC en agricultura y desarrollo rural: apoyo para artículos introductorios, de contexto, marco conceptual. Repositorio CEPAL \\
\midrule
Revista Colombiana (2025) & Retos de la Agricultura Digital & Estudia barreras de digitalización rural en Colombia (infraestructura, conectividad, desigualdades): útil para análisis crítico de adopción tecnológica. Dialnet \\
\midrule
Estudio (2024) & How do coffee farmers engage with digital technologies? A capabilities perspective & Analiza la adopción tecnológica en caficultores colombianos desde la perspectiva de capacidades y contexto local. SpringerLink \\
\midrule
Artículo (2025) & Digital integration toward smart and sustainable agriculture in the European region & Examina impactos de la digitalización en sostenibilidad agrícola. ScienceDirect \\
\midrule
Estudio de caso (2025) & Smart Rural Communities: Action Research in Colombia and Mozambique & Explora características de comunidades rurales que adoptan soluciones digitales agrícolas. MDPI \\
\midrule
Artículo (2024) & Blockchain and agricultural sustainability in South America: a systematic review & Revisión del uso de blockchain en agricultura sostenible en Sudamérica. Frontiers \\
\midrule
Artículo (2024) & Los mercados campesinos como estrategia de sustentabilidad en los sistemas alimentarios de Cundinamarca, Colombia & Analiza mercados campesinos y su contribución a sostenibilidad, equidad y desarrollo territorial. Hemeroteca UNAD \\
\midrule
Artículo (2022) & ¿Sentando las bases para una sociedad más consciente? Cómo vendedores, consumidores y organizadores construyen socialmente los mercados campesinos en Bogotá, Colombia & Describe dinámicas de mercado campesino, relaciones sociales, organización comunitaria y su rol en seguridad alimentaria urbana. Redalyc \\
\midrule
IICA (2025) & Digitalización de la asistencia técnica para integrar la agricultura familiar a cadenas de valor & Buen respaldo institucional actual sobre importancia de TIC para asistencia técnica, formación y acceso al mercado. iica.int \\
\midrule
FAO (2023) & 1000 Digital Villages in Latin America and the Caribbean & Describe casos de comercialización directa desde finca al consumidor, uso de plataformas digitales, fortalecimiento rural. FAOHome \\
\bottomrule
\end{tabular}
}
\caption{Referencias adicionales para la investigación sobre adopción tecnológica en el agro colombiano.}
\label{tab:referencias_adicionales}
\end{table}

\section{Detalles Adicionales sobre Economía Circular}

La economía circular es un modelo que busca eliminar el desperdicio y promover el uso continuo de recursos. En el contexto agrícola colombiano, esto implica aprovechar subproductos como cáscaras, bagazo y residuos orgánicos para generar nuevos ingresos.

Por ejemplo, las cáscaras de frutas pueden ser utilizadas para compostaje, alimentación animal o incluso artesanías. Esto no solo reduce el impacto ambiental, sino que crea oportunidades económicas para los productores rurales.

La implementación en el portal incluye una sección dedicada a la venta de estos subproductos, conectando productores con compradores interesados en materiales sostenibles. Además, se ofrecen capacitaciones sobre técnicas de economía circular para maximizar los beneficios.

Este enfoque contribuye a la sostenibilidad a largo plazo, fortaleciendo la resiliencia de las comunidades agrícolas frente a cambios climáticos y económicos.

\section{Impacto del Cambio Climático en la Agricultura}

El cambio climático representa una amenaza significativa para la agricultura colombiana, con patrones de lluvia irregulares, temperaturas extremas y aumento de plagas. Los productores necesitan herramientas para adaptarse, como pronósticos climáticos precisos y recomendaciones de cultivos resistentes.

El portal integra datos climáticos históricos y predicciones para ayudar a los agricultores a planificar sus siembras. Por ejemplo, alertas de heladas permiten proteger cultivos vulnerables, reduciendo pérdidas.

Además, se promueve la diversificación de cultivos para mitigar riesgos. Esta adaptación no solo preserva la producción, sino que asegura la seguridad alimentaria en regiones afectadas.

\section{Marketing Digital para Productores}

El marketing digital permite a los productores llegar a nuevos mercados urbanos. Herramientas como publicaciones en redes sociales y descripciones atractivas de productos elevan el valor percibido.

El portal ofrece plantillas para crear contenido visual, destacando la historia y calidad de los productos locales. Esto fomenta el comercio directo y fortalece la identidad cultural de las comunidades rurales.

Con el marketing digital, los productores pueden competir en un mercado global, aumentando sus ingresos y visibilidad.

\section{Blockchain en la Trazabilidad Agrícola}

El Blockchain es una tecnología que registra datos de forma inalterable y transparente. En la agricultura, garantiza la autenticidad de los productos, facilitando exportaciones y confianza del consumidor.

Implementación en el proyecto: El portal generará certificados digitales de trazabilidad para cada lote, con códigos QR únicos que verifican el historial del producto.

\section{Pagos Digitales y Financiación}

La digitalización de pagos elimina efectivo y demoras. Billeteras como Nequi y Daviplata facilitan transacciones seguras y acceso a financiamiento.

El portal integra estas opciones, permitiendo pagos directos y promoviendo inclusión financiera en áreas rurales.

\section{Cooperación entre Agricultores}

La economía colaborativa permite compras grupales y envíos conjuntos, reduciendo costos y aumentando competitividad.

El portal facilita asociaciones virtuales, fortaleciendo lazos comunitarios y mejorando el acceso a mercados.

\section{Soberanía Alimentaria}

La soberanía alimentaria implica control local sobre producción y distribución. El portal promueve semillas nativas y prácticas agroecológicas.

Esto asegura alimentos sanos y estabilidad económica, protegiendo el patrimonio agrícola regional.

\section{Visión Futurista del Agro}

El futuro incluye sensores IoT, IA y riego automatizado. El portal sienta bases para estas innovaciones, transformando al agricultor en empresario tecnológico.

Colombia puede liderar en agricultura 5.0 con inversión en conectividad y educación.

\section{Logística y Cadena de Suministro}

La logística eficiente es crucial para reducir pérdidas post-cosecha. El portal optimiza rutas y coordina transporte compartido.

Implementación: Módulos para seguimiento en tiempo real y predicción de demoras, mejorando entrega y calidad.

\section{Análisis de Mercado}

Tableros analíticos proporcionan insights sobre tendencias de precios y demanda. Esto permite decisiones informadas para maximizar ganancias.

El portal integra datos de mercado para recomendaciones personalizadas.

\section{Participación Comunitaria}

El proyecto involucra comunidades en diseño y pruebas. Esto asegura que el portal responda a necesidades reales.

Talleres participativos fortalecen apropiación y sostenibilidad.

\section{Impacto Ambiental}

Prácticas sostenibles como economía circular reducen huella ecológica. El portal promueve agricultura regenerativa.

Esto contribuye a resiliencia climática y biodiversidad.

\section{Integración Tecnológica}

El portal integra IA para predicciones de cosecha y optimización de recursos. Sensores IoT monitorean condiciones del suelo y clima.

Esto permite agricultura de precisión, maximizando productividad y minimizando insumos.

\section{Educación Continua}

Programas de formación digital incluyen cursos en línea y talleres. El portal ofrece certificaciones para motivar aprendizaje.

Esto asegura que los productores mantengan habilidades actualizadas.

\section{Redes de Colaboración}

El proyecto fomenta alianzas entre productores, instituciones y empresas. Redes colaborativas comparten conocimientos y recursos.

Esto fortalece el ecosistema agrícola regional.

\section{Evaluación de Impacto}

Métricas de éxito incluyen aumento en ingresos, reducción de pérdidas y adopción tecnológica. Evaluaciones periódicas guían mejoras.

El portal incluye herramientas para autoevaluación y retroalimentación.

\section{Casos de Estudio}

Ejemplos reales de productores que han adoptado el portal muestran beneficios tangibles. Historias de éxito inspiran a otros.

Esto demuestra viabilidad y potencial transformador.

\section{Desafíos Futuros}

Identificación de obstáculos como resistencia cultural y brechas digitales. Estrategias para superarlos incluyen alianzas y innovación.

El proyecto se adapta continuamente a cambios.

\section{Conclusiones del Apéndice}

El portal representa un avance significativo en digitalización agrícola. Con compromiso, puede transformar el sector en Colombia.

Recomendaciones incluyen expansión y monitoreo continuo.